% Figure for Thermal Physics §A2.1 (van der Waals isotherms in reduced units p/p_c against V/V_c, with an ideal-gas isotherm and the liquid-gas plateau; after University Physics Vol. 2, §2.1).
\begin{tikzpicture}[font=\small]
  \begin{axis}[nfig, width=9.5cm, height=6.5cm, xmin=0.3, xmax=4.2, ymin=0, ymax=2.4,
      xlabel={$V/V_c$}, ylabel={$p/p_c$},
      xtick={1,2,3,4}, ytick={0.5,1,1.5,2},
      xlabel style={at={(axis description cs:1.02,0)}, anchor=west},
      ylabel style={rotate=-90, at={(axis description cs:0,1)}, anchor=south}, clip=true, restrict y to domain=-1:2.6, unbounded coords=jump]
    \addplot[figR, domain=0.505:4.0, samples=160] {8*1.5/(3*x-1)-3/x^2};
    \addplot[figR, domain=0.47:4.0, samples=160] {8*1.2/(3*x-1)-3/x^2};
    \addplot[black, domain=0.452:4.0, samples=200] {8*1.0/(3*x-1)-3/x^2};
    \addplot[figB, domain=0.445:4.0, samples=200] {8*0.9/(3*x-1)-3/x^2};
    \addplot[figB, very thick, domain=0.6034:2.3488] {0.6470};
    \addplot[figR, dashed, thin, domain=1.12:4.0, samples=60] {(8/3)*1.5/x};
    \addplot[only marks, mark=*, mark size=1.5pt, black] coordinates {(1,1)};
    \node[figR, font=\scriptsize, anchor=south] at (axis cs:3.75,1.06) {$1.5\,T_c$};
    \node[figR, font=\scriptsize, anchor=south] at (axis cs:3.62,0.83) {$1.2\,T_c$};
    \node[font=\scriptsize, anchor=south] at (axis cs:3.62,0.635) {$T_c$};
    \node[figB, font=\scriptsize, anchor=north] at (axis cs:3.75,0.465) {$0.9\,T_c$};
    \node[figR, font=\scriptsize, anchor=west] at (axis cs:2.05,2.05) {ideal gas, $1.5\,T_c$ (dashed)};
    \node[font=\scriptsize, anchor=south west] at (axis cs:1.05,1.2) {critical point};
    \draw[black!55, thin] (axis cs:1.1,1.21) -- (axis cs:1.02,1.04);
    \node[figB, font=\scriptsize] at (axis cs:1.72,0.565) {liquid + vapour};
    \node[figB, font=\scriptsize, anchor=west] at (axis cs:0.75,0.25) {oscillation (unphysical)};
    \draw[black!55, thin] (axis cs:0.85,0.3) -- (axis cs:0.74,0.41);
  \end{axis}
\end{tikzpicture}
